\section{Effective antipodal expansion and stable sections}
\label{sec:effective52}
The one-surplus constant admits an explicit lower bound. The first
step is a trace constraint for a stochastic extension of an antipodal
map; the second converts support displacement into a volume increment.
For a convex body containing zero in its interior write
\[
 s_0(P)=\min\{\lambda\ge1:-P\subset\lambda P\}.
\]
This is asymmetry about the specified origin, not minimization over
translations.

\begin{lemma}[A trace bound for polyhedral asymmetry]\label{lem:trace52}
Let $P\subset\R^D$ be a full-dimensional polytope containing zero
in its interior. If $P$ has at most $k$ vertices or at most $k$ facets,
where $D+1\le k<2D$, then
\begin{equation}\label{eq:trace-asymmetry52}
                          s_0(P)\ge\frac{D}{k-D}.
\end{equation}
\end{lemma}
\begin{proof}
First suppose $P$ has $m\le k$ vertices, written as the rows of
$Z\in\R^{m\times D}$. For $\lambda=s_0(P)$, choose a convex
representation of each $-Z(i,\cdot)/\lambda$ in the vertices.
The resulting nonnegative row-stochastic matrix $T$ satisfies
\[
              T\mathbf1=\mathbf1,\qquad TZ=-\lambda^{-1}Z.
\]
The columns of $[\mathbf1,Z]$ are independent. Thus $1$ is an
eigenvalue of $T$, and $-\lambda^{-1}$ has geometric multiplicity
at least $D$. All eigenvalues of a row-stochastic matrix have modulus
at most one. The remaining $m-D-1$ eigenvalues, counted algebraically,
therefore have real parts summing to at most $m-D-1$. Since the
trace is nonnegative,
\[
           0\le\operatorname{tr}T\le m-D-D/\lambda.
\]
This implies $\lambda\ge D/(m-D)\ge D/(k-D)$.

For the facet case use the polar body $P^\circ$. It has at most
$k$ vertices. Polarity reverses inclusions and commutes with negation;
indeed $-P\subset\lambda P$ is equivalent to
$\lambda^{-1}P^\circ\subset -P^\circ$, hence to
$-P^\circ\subset\lambda P^\circ$. Thus
$s_0(P)=s_0(P^\circ)$, and the first part applies.
\end{proof}

\begin{lemma}[From support displacement to volume]\label{lem:cap-volume52}
Suppose $r\Ball^D\subset P\subset[-1,1]^D$ and $s_0(P)\ge L>1$.
Let $\omega_{D-1}$ denote the volume of the unit ball in $\R^{D-1}$.
Then
\begin{equation}\label{eq:cap-volume52}
 \frac{\operatorname{vol}(\conv(P\cup(-P)))}{\operatorname{vol}(P)}
 \ge1+\frac{\omega_{D-1}r^D}{D2^D}
                    (L-1)(1-L^{-1})^{D-1}.
\end{equation}
\end{lemma}
\begin{proof}
The maximum of $h_P(-u)/h_P(u)$ over unit $u$ is $s_0(P)$, with
positive denominator. Choose a maximizing $u$, write $h=h_P(u)\ge r$,
and choose $x\in-P$ with $t=\ip u x=h_P(-u)\ge Lh$.
The disk $\{z:\ip u z=0,\ \|z\|\le r\}$ lies in $P$.
Its cone with apex $x$ lies in $\conv(P\cup(-P))$.
The portion with $\ip u z>h$ lies outside $P$. Its height is
$t-h\ge(L-1)r$, and its base in the plane $\ip u z=h$ is a
translated disk of radius $r(1-h/t)\ge r(1-L^{-1})$.
The cone-volume formula gives an added volume at least
$\omega_{D-1}r^D(L-1)(1-L^{-1})^{D-1}/D$.
Dividing by $\operatorname{vol}(P)\le2^D$ proves the claim.
The cone may be oblique; the formula uses perpendicular height and
therefore is unchanged.
\end{proof}

\begin{theorem}[An effective one-surplus multiplier]\label{thm:beta52}
Whenever $\mathscr H_{D,\rho}$ is nonempty and $D\ge3$,
\begin{equation}\label{eq:beta52}
 \beta_D(\rho)\ge
 1+\frac{\omega_{D-1}(D-2)^D}{2^{D+1}D^{3D/2}}\rho^D.
\end{equation}
In dimension three, the simpler rational bound is
\begin{equation}\label{eq:beta3-rational52}
                         \beta_3(\rho)\ge1+\rho^3/1024.
\end{equation}
Consequently, for every signed-permutation alphabet in dimension
three containing $I,-I$,
\begin{equation}\label{eq:six-effective52}
 (1+\rho^3/1024)^N>6\rho^{-3}
                    \quad\Longrightarrow\quad W_{N,0}=6.
\end{equation}
For $\rho=1/10$, a sufficient integer horizon is
$N_0=9\,216\,009$. Neither this horizon nor the multiplier is claimed
optimal.
\end{theorem}
\begin{proof}
Each member of $\mathscr H_{D,\rho}$ contains the Euclidean ball
of radius $r=\rho/\sqrt D$. Lemma~\ref{lem:trace52}, with $k=D+2$,
gives $L=D/2$. Substitution in Lemma~\ref{lem:cap-volume52} gives
\eqref{eq:beta52} for every member, hence for the minimum.
For $D=3$ its coefficient is $\pi/(1296\sqrt3)>1/864>1/1024$,
using $\pi>3$ and $\sqrt3<2$. Apply the occupation inequality
\eqref{eq:surplus-budget51}; the six-state permutation construction
supplies the upper bound.

For the numerical horizon put $x=1/1\,024\,000$. The elementary
inequality $\log(1+x)\ge x/(1+x)$ gives
$N_0\log(1+x)\ge9$. The finite rational Taylor sum
$\sum_{j=0}^{20}9^j/j!>6000$ proves $e^9>6000=6\rho^{-3}$.
Thus the strict inequality in \eqref{eq:six-effective52} holds.
\end{proof}

\subsection{Quantitative robustness beyond minimum rank}
The unconditional noisy classification in Section~\ref{sec:noisy52}
requires no conditioning. A separate conditioning quantity recovers
explicit geometric multipliers, including a noisy one-surplus bound.
It is specified as data, not assumed automatically uniform in time.

At a command cut let $E_t$ be a matrix whose $r_t$ rows are linearly
independent actual reachable probability rows and span all reachable
rows. Put $P_t=\operatorname{rowspan}(E_t)\cap\Delta_{K_t-1}$.
Every $p\in P_t$ has a unique expression $p=\alpha E_t$. Define
\begin{equation}\label{eq:conditioning52}
                 A_t=\max_{\alpha E_t\in P_t}\|\alpha\|_1.
\end{equation}
The maximum exists: the coefficient map is a continuous linear
inverse on the reachable span, and $P_t$ is compact. Also $A_t\ge1$.
Different reachable bases can give different values; a stated bound
means that such a basis has been supplied or shown to exist.

\begin{theorem}[Stable physical sections at arbitrary width]
\label{thm:conditioned52}
Let the machine have worst-word mean error $\delta=2\varepsilon$
with $0\le\delta<\rho/(2D)$. At selected cuts suppose bases as above
satisfy $A_t\le A$. Let $Z_t$ be the actual identity-suffix mean
matrix and set $S_t=\{(pZ_t)^T:p\in P_t\}$. Then
\[
 C_b\subset S_t\subset[-1,1]^D,\quad
 f_0(S_t)\le\binom{K_t}{r_t-1},\quad b=\rho-D\delta>0.
\]
The affine observation rank is $D+1$ including mass, and the fiber
dimension is $r_t-D-1$. For selected $s<t$ and every intervening word,
\begin{equation}\label{eq:conditioned-containment52}
 U_w S_s\subset S_t+h[-1,1]^D\subset\Lambda S_t,
 \qquad h=A(1+\sqrt D)\delta,\quad \Lambda=1+Dh/b.
\end{equation}
Thus, for $m\ge1$ selected cuts of width at most $K$,
\begin{equation}\label{eq:conditioned-budget52}
 \left(\frac{\gamma_{\calA}(V_K,b)}{\Lambda^D}\right)^{m-1}
                                  \le D!b^{-D}.
\end{equation}
When $I,-I\in\calA$, the same bound with $\beta_D(b)$ in place of
$\gamma_{\calA}(V_K,b)$ holds for selected cuts of width at most
$D+2$, for $D\ge3$.
\end{theorem}
\begin{proof}
The affine-hull and vertex arguments of Theorem~\ref{thm:section51}
use only reachability and stochasticity and therefore remain valid.
The signed identity-prefix means are within $\delta$ per coordinate
of $\pm\rho e_i$. Their support functions imply $C_b\subset S_t$
exactly as in Lemma~\ref{lem:anchor51}. Full-dimensionality then gives
rank $D+1$ for the mass-and-observation map, proving the fiber formula.
This also recovers the one-surplus vertex/facet dichotomy.

For any suffix $v$, the target mean following each actual basis prefix
is the target identity-suffix mean multiplied by $U_v^T$.
Consequently every entry of
$E_s(B_s(v)-Z_sU_v^T)$ has modulus at most $(1+\sqrt D)\delta$.
Multiplying by the coefficient row $\alpha$ of $p\in P_s$ bounds
its observable defect by $h$. Reachable spans are invariant under
actual stochastic transitions even without exactness, so
$P_sT_w\subset P_t$. Take $v$ to be $w$ followed by the full identity
suffix from $t$ to the terminal cut; then $B_s(v)=T_wZ_t$.
This proves the first containment in \eqref{eq:conditioned-containment52}.
The inclusion $h[-1,1]^D\subset(Dh/b)S_t$ gives the second.

Identity padding shows that the next selected body, dilated by
$\Lambda$, contains every generator image of the preceding body.
Use the definition of $\gamma_{\calA}(V_K,b)$ and iterate volume
ratios. For one surplus label, use both signs and the vertex/facet
dichotomy instead. The endpoint volumes are bounded below by
$2^Db^D/D!$ and above by $2^D$, giving both assertions.
\end{proof}

\begin{corollary}[An explicit conditioned noisy six-state interval]
\label{cor:conditioned-six52}
In dimension three let $I,-I\in\calA$, with all commands signed
permutations. Fix $A\ge1$. Among machines admitting reachable bases
with $A_t\le A$ at every cut, the optimal peak is six whenever
\begin{equation}\label{eq:conditioned-six52}
 0\le\varepsilon\le\frac{\rho^4}{2^{23}A},\qquad
 (1+\rho^3/16384)^N>48\rho^{-3}.
\end{equation}
The error bound is independent of the horizon. It is an explicit
statement about this conditioned class, not a claim that arbitrary
machines have uniformly conditioned reachable bases.
\end{corollary}
\begin{proof}
Write $\delta=2\varepsilon\le\rho^4/(2^{22}A)$. Then $b\ge\rho/2$
and, since $1+\sqrt3<3$, $u=\Lambda-1\le18A\delta/\rho$.
Here $u<1/4$, so $(1+u)^3\le1+4u\le1+\rho^3/32768$.
By \eqref{eq:beta3-rational52}, $\beta_3(b)\ge1+\rho^3/8192$.
It follows that
\[
 \frac{\beta_3(b)}{\Lambda^3}
 \ge\frac{1+\rho^3/8192}{1+\rho^3/32768}
 \ge1+\rho^3/16384.
\]
The conditioned occupation bound excludes peak at most five under
\eqref{eq:conditioned-six52}. The exact six-label machine has all
six point masses reachable at every cut by identity prefixes. Choose
these as its basis, giving $A_t=1$. It attains the upper bound inside
the conditioned class.
\end{proof}

\subsection{An unconditional computable noise interval}
There is also an unconditional positive-error six-state theorem. Its
constant is computable from finite algebraic data, though the following
algorithm is not a practical way to obtain a sharp value.

\begin{theorem}[Effective finite-horizon separation propagated to all horizons]
\label{thm:computable-six52}
Let $D=3$, let $\rho\in(0,1)$ be real algebraic, and let $\calA$ be
any signed-permutation alphabet containing $I,-I$. Given any integer
$N_0$ satisfying \eqref{eq:six-effective52}, there is an algorithm
which terminates with a positive dyadic rational $\varepsilon_0$ such
that
\[
          W_{N,\varepsilon}=6\quad
          (N\ge N_0,\ 0\le\varepsilon\le\varepsilon_0).
\]
This assertion has no rank or conditioning restriction.
\end{theorem}
\begin{proof}
Pad every register of a peak-at-most-five machine at horizon $N_0$
to five labels. Its initialization, command rows, and decoder means
belong to a compact product of simplexes and intervals. The maximum
terminal error over the finitely many words, seeds, and queries is
continuous, so its minimum $\eta$ is attained. By
Theorem~\ref{thm:beta52}, no exact such machine exists; hence $\eta>0$.

Feasibility at any algebraic tolerance is decidable over the real
algebraic numbers. One explicit formula introduces actual prefix
probability rows $q_{x,u}$, the common epoch-and-letter matrices
$T_{t,a}$, and bounded decoder columns $d_j$. Impose initialization,
$q_{x,ua}=q_{x,u}T_{|u|+1,a}$, and
\[
 |q_{x,w}d_j-\rho e_j^TU_wx|\le2\varepsilon
 \quad (|w|=N_0).
\]
All variables lie in a compact box, stochasticity is explicit, and
the common rows enforce compatibility. Enumerating words may require
exponentially many variables; no polynomial complexity is claimed.
Algebraic coefficients are represented by defining polynomials and
rational isolating intervals. Real quantifier elimination decides the
closed feasibility formula \cite{BPR}.

Test $\varepsilon=2^{-j}$ successively for $j=1,2,\ldots$, stopping
at the first infeasible value. Since $2^{-j}<\eta$ eventually, this
terminates and returns $\varepsilon_0$. For a longer horizon, final
identity padding and composition into the decoder would turn any
peak-five machine into one at $N_0$ with no larger error, a
contradiction. The exact six-state machine proves equality.
\end{proof}

The theorem specifies a terminating computation, not its execution on
the large illustrative horizon $9\,216\,009$. No numerical value of
the unconditional $\varepsilon_0$ at that horizon is reported. The
explicit interval \eqref{eq:conditioned-six52} and the unrestricted
computable interval are different conclusions and should not be
interchanged.
